\chapter{Về Bạn Bè}
\markboth{Về Bạn Bè}{About Friendship}

\section{Chọn Bạn Tốt}
\begin{truyen}{Chọn Bạn Tốt}{Choose Good Friends}
\chuhoa{B}{ạn bè ảnh hưởng đến cách nghĩ, thói quen và cả hướng đi của em nhiều hơn em tưởng.}
\textit{(Friends influence your thinking, habits, and direction more than you may realize.)}

Chọn bạn tốt không có nghĩa là tìm người hoàn hảo.
\textit{(Choosing good friends does not mean finding perfect people.)}

Nó là chọn những người làm mình muốn tử tế hơn, chăm hơn và sáng hơn.
\textit{(It means choosing people who make you want to be kinder, more diligent, and brighter.)}

Tình bạn đúng giúp mình lớn lên đúng hơn.
\textit{(Right friendship helps us grow better.)}
\end{truyen}

\begin{baihoc}
Bạn bè không chỉ đi cùng mình; họ còn âm thầm kéo mình theo một hướng nào đó.
\textit{(Friends do not only walk beside us; they also quietly pull us in some direction.)}
\end{baihoc}

\begin{ghinhoanh}
\textbf{Short English to remember:}

Good friends shape good growth.

\end{ghinhoanh}

\begin{tuhocnhanh}
1. Vì sao cần chọn bạn tốt? \textit{Why is it important to choose good friends?}

2. Bạn tốt không có nghĩa là gì? \textit{What does a good friend not mean?}

3. Bạn bè đang kéo em theo hướng nào? \textit{In what direction are your friends pulling you?}
\end{tuhocnhanh}
\ngancach

\section{Tôn Trọng Bạn}
\begin{truyen}{Tôn Trọng Bạn}{Respect Your Friends}
\chuhoa{T}{ôn trọng là nền của tình bạn đẹp.}
\textit{(Respect is the foundation of good friendship.)}

Nó thể hiện ở cách lắng nghe, cách nói chuyện, và cách mình không coi thường cảm xúc hay khác biệt của bạn.
\textit{(It appears in listening, speaking, and refusing to belittle a friend's feelings or differences.)}

Thiếu tôn trọng, tình bạn rất dễ thành nơi làm nhau mệt.
\textit{(Without respect, friendship easily becomes tiring.)}

Biết tôn trọng bạn là một dấu hiệu của sự trưởng thành.
\textit{(Knowing how to respect friends is a sign of maturity.)}
\end{truyen}

\begin{baihoc}
Tình bạn bền không chỉ nhờ vui vẻ, mà còn nhờ sự tôn trọng nhau.
\textit{(Lasting friendship depends not only on fun, but also on respect.)}
\end{baihoc}

\begin{ghinhoanh}
\textbf{Short English to remember:}

Respect keeps friendship healthy.

\end{ghinhoanh}

\begin{tuhocnhanh}
1. Vì sao tôn trọng là nền của tình bạn? \textit{Why is respect the base of friendship?}

2. Nó thể hiện qua những gì? \textit{How does it show itself?}

3. Em có đang tôn trọng bạn bè đủ không? \textit{Are you respecting your friends enough?}
\end{tuhocnhanh}
\ngancach

\section{Giúp Bạn Khi Cần}
\begin{truyen}{Giúp Bạn Khi Cần}{Help Friends When Needed}
\chuhoa{G}{iúp bạn đúng lúc là một điều nhỏ nhưng rất ấm.}
\textit{(Helping a friend at the right time is small but warm.)}

Có khi chỉ là giảng lại bài, nhắc việc, hoặc ngồi cạnh một lúc khi bạn đang rối.
\textit{(Sometimes it is simply re-explaining a lesson, reminding a task, or sitting beside a struggling friend.)}

Sự giúp đỡ ấy làm tình bạn bớt nông và bớt chỉ là đi chơi cùng nhau.
\textit{(That help makes friendship deeper than merely having fun together.)}

Người biết giúp bạn thường cũng đang học cách sống tử tế hơn.
\textit{(Those who help friends are also learning to live more kindly.)}
\end{truyen}

\begin{baihoc}
Giúp một người lúc cần là một cách làm cho tình bạn có trọng lượng hơn.
\textit{(Helping someone when needed gives friendship more real weight.)}
\end{baihoc}

\begin{ghinhoanh}
\textbf{Short English to remember:}

Timely help strengthens friendship.

\end{ghinhoanh}

\begin{tuhocnhanh}
1. Vì sao nên giúp bạn khi cần? \textit{Why should you help friends when needed?}

2. Sự giúp đỡ làm tình bạn khác đi thế nào? \textit{How does help change friendship?}

3. Em có thể giúp bạn điều gì ngay lúc này? \textit{What can you help a friend with now?}
\end{tuhocnhanh}
\ngancach

\section{Không Cười Khi Bạn Sai}
\begin{truyen}{Không Cười Khi Bạn Sai}{Do Not Laugh When a Friend Is Wrong}
\chuhoa{S}{ai trong lúc học là điều rất bình thường, nên đem nó ra làm trò cười dễ làm bạn mình co lại.}
\textit{(Being wrong while learning is normal, so turning it into a joke can make a friend shrink back.)}

Một tiếng cười thiếu tử tế có thể làm ai đó sợ hỏi, sợ phát biểu và sợ thử lại.
\textit{(An unkind laugh can make someone afraid to ask, speak, or try again.)}

Không cười khi bạn sai là một cách giữ cho lớp học và tình bạn an toàn hơn.
\textit{(Not laughing when a friend is wrong keeps class and friendship safer.)}

Đó cũng là một nét đẹp của sự trưởng thành.
\textit{(It is also a sign of maturity.)}
\end{truyen}

\begin{baihoc}
Tử tế với lỗi sai của người khác là một cách rất thật để làm bạn.
\textit{(Being kind toward another person's mistakes is a real way of being a friend.)}
\end{baihoc}

\begin{ghinhoanh}
\textbf{Short English to remember:}

Kindness matters when others fail.

\end{ghinhoanh}

\begin{tuhocnhanh}
1. Vì sao không nên cười khi bạn sai? \textit{Why should you not laugh when a friend is wrong?}

2. Một tiếng cười thiếu tử tế có thể gây gì? \textit{What can an unkind laugh cause?}

3. Em có đang đủ dịu với lỗi của bạn không? \textit{Are you gentle enough with your friends' mistakes?}
\end{tuhocnhanh}
\ngancach

\section{Không Ganh Tị}
\begin{truyen}{Không Ganh Tị}{Do Not Be Jealous}
\chuhoa{G}{anh tị làm con người mệt và làm tình bạn bị đục.}
\textit{(Jealousy exhausts people and clouds friendship.)}

Khi nhìn thấy bạn giỏi hơn ở một điểm nào đó, em có thể chọn học từ bạn thay vì khó chịu với bạn.
\textit{(When you see a friend stronger in something, you can choose to learn from them instead of resenting them.)}

Không ganh tị không có nghĩa là không buồn chút nào.
\textit{(Not being jealous does not mean never feeling a sting.)}

Nó nghĩa là không để cảm xúc ấy dẫn mình đi sai.
\textit{(It means not letting that feeling lead you in the wrong direction.)}
\end{truyen}

\begin{baihoc}
Nhìn cái tốt của bạn bằng con mắt học hỏi thường tốt hơn nhìn bằng ganh tị.
\textit{(Seeing a friend's strength with a learning eye is better than seeing it with envy.)}
\end{baihoc}

\begin{ghinhoanh}
\textbf{Short English to remember:}

Learn from strength; do not envy it.

\end{ghinhoanh}

\begin{tuhocnhanh}
1. Vì sao ganh tị có hại cho tình bạn? \textit{Why is jealousy harmful to friendship?}

2. Em có thể thay ganh tị bằng điều gì? \textit{What can you replace jealousy with?}

3. Em đang học được gì từ điểm mạnh của bạn? \textit{What are you learning from a friend's strength?}
\end{tuhocnhanh}
\ngancach

\section{Học Từ Bạn}
\begin{truyen}{Học Từ Bạn}{Learn from Friends}
\chuhoa{B}{ạn bè không chỉ để đi cùng, mà còn có thể là những người dạy mình rất nhiều.}
\textit{(Friends are not only companions; they can also teach us many things.)}

Có bạn cẩn thận hơn, có bạn bình tĩnh hơn, có bạn chăm hơn hoặc sáng tạo hơn.
\textit{(Some friends are more careful, calmer, harder-working, or more creative.)}

Biết học từ bạn giúp em bớt tự ái và lớn lên nhanh hơn.
\textit{(Learning from friends helps you become less defensive and grow faster.)}

Một tình bạn đẹp thường có cả sự học hỏi lẫn nhau.
\textit{(A good friendship often contains mutual learning.)}
\end{truyen}

\begin{baihoc}
Bạn bè tốt không chỉ làm mình vui; họ còn làm mình khá hơn.
\textit{(Good friends not only make us happy; they also help us improve.)}
\end{baihoc}

\begin{ghinhoanh}
\textbf{Short English to remember:}

Friends can teach in quiet ways.

\end{ghinhoanh}

\begin{tuhocnhanh}
1. Vì sao nên học từ bạn bè? \textit{Why should you learn from friends?}

2. Bạn bè có thể dạy mình những gì? \textit{What can friends teach us?}

3. Em học được điều gì từ một người bạn? \textit{What have you learned from a friend?}
\end{tuhocnhanh}
\ngancach

\section{Chia Sẻ Kiến Thức}
\begin{truyen}{Chia Sẻ Kiến Thức}{Share Knowledge}
\chuhoa{C}{hỉ giữ phần mình biết cho riêng mình không làm lớp học đẹp hơn.}
\textit{(Keeping what you know only to yourself does not make the classroom better.)}

Khi chia sẻ kiến thức đúng mực, em đang giúp bạn đỡ lạc và cũng làm mình hiểu chắc hơn.
\textit{(When you share knowledge properly, you help a friend get less lost and strengthen your own understanding.)}

Điều này không làm em mất đi lợi thế thật sự.
\textit{(This does not take away your real advantage.)}

Nó chỉ làm tri thức trở nên có ích hơn cho nhiều người.
\textit{(It simply makes knowledge more useful to more people.)}
\end{truyen}

\begin{baihoc}
Điều mình biết, khi chia sẻ đúng cách, thường không mất đi mà còn sáng rõ hơn.
\textit{(What we know, when shared well, is often not lost but clarified.)}
\end{baihoc}

\begin{ghinhoanh}
\textbf{Short English to remember:}

Shared knowledge can strengthen everyone.

\end{ghinhoanh}

\begin{tuhocnhanh}
1. Vì sao nên chia sẻ kiến thức? \textit{Why should you share knowledge?}

2. Nó giúp chính em như thế nào? \textit{How does it help you too?}

3. Em có thể chia sẻ gì với bạn hôm nay? \textit{What can you share with a friend today?}
\end{tuhocnhanh}
\ngancach

\section{Khích Lệ Nhau}
\begin{truyen}{Khích Lệ Nhau}{Encourage One Another}
\chuhoa{M}{ột lời động viên đúng lúc đôi khi giữ bạn mình ở lại với việc học thêm một đoạn nữa.}
\textit{(A timely word of encouragement can keep a friend with learning a little longer.)}

Tuổi học trò có rất nhiều lúc nản, tự ti và muốn bỏ.
\textit{(Student life has many moments of discouragement, insecurity, and the desire to quit.)}

Khi bạn bè biết khích lệ nhau, lớp học có thêm một nguồn sức mạnh rất tử tế.
\textit{(When friends know how to encourage each other, the classroom gains a kind source of strength.)}

Khích lệ nhau là một phần đẹp của tình bạn.
\textit{(Encouragement is a beautiful part of friendship.)}
\end{truyen}

\begin{baihoc}
Không phải lúc nào bạn bè cũng cần lời khuyên dài; đôi khi chỉ cần một lời khích lệ đúng lúc.
\textit{(Friends do not always need long advice; sometimes they only need one timely encouraging sentence.)}
\end{baihoc}

\begin{ghinhoanh}
\textbf{Short English to remember:}

Encouragement can keep effort alive.

\end{ghinhoanh}

\begin{tuhocnhanh}
1. Vì sao nên khích lệ nhau? \textit{Why should friends encourage each other?}

2. Lời khích lệ có thể giữ lại điều gì? \textit{What can encouragement preserve?}

3. Em có thể nói lời nào để nâng một người bạn lên? \textit{What words can lift a friend today?}
\end{tuhocnhanh}
\ngancach

\section{Cạnh Tranh Lành Mạnh}
\begin{truyen}{Cạnh Tranh Lành Mạnh}{Healthy Competition}
\chuhoa{C}{ạnh tranh không xấu nếu nó giúp mình cố gắng hơn mà không làm mình xấu đi.}
\textit{(Competition is not bad if it pushes us forward without making us worse.)}

Cạnh tranh lành mạnh là cùng nhau tiến lên, không phải tìm cách dìm nhau xuống.
\textit{(Healthy competition means moving upward together, not pulling each other down.)}

Nó làm em nhìn bạn như một động lực, không phải một mối đe dọa.
\textit{(It lets you see a friend as motivation, not as a threat.)}

Đây là một thái độ rất nên có trong tuổi học trò.
\textit{(This is a very good attitude for student life.)}
\end{truyen}

\begin{baihoc}
Cạnh tranh đẹp là cạnh tranh làm cả hai bên tốt hơn.
\textit{(Good competition is competition that makes both sides better.)}
\end{baihoc}

\begin{ghinhoanh}
\textbf{Short English to remember:}

Healthy competition improves without harming.

\end{ghinhoanh}

\begin{tuhocnhanh}
1. Vì sao cạnh tranh không hẳn là xấu? \textit{Why is competition not always bad?}

2. Cạnh tranh lành mạnh khác gì với ganh đua xấu? \textit{How is healthy competition different from harmful rivalry?}

3. Em có đang cạnh tranh theo cách đẹp không? \textit{Are you competing in a healthy way?}
\end{tuhocnhanh}
\ngancach

\section{Tình Bạn Đẹp}
\begin{truyen}{Tình Bạn Đẹp}{Beautiful Friendship}
\chuhoa{T}{ình bạn đẹp không cần lúc nào cũng ồn ào hay hoàn hảo.}
\textit{(A beautiful friendship does not need to be loud or perfect.)}

Có khi nó chỉ là sự ở cạnh tử tế, sự tôn trọng, sự giúp nhau đúng lúc, và cảm giác an toàn khi ở bên nhau.
\textit{(Sometimes it is simply kind presence, respect, timely help, and the feeling of safety with each other.)}

Tuổi học trò có một vài tình bạn đẹp là một điều rất quý.
\textit{(Having a few beautiful friendships in student life is deeply precious.)}

Nó làm những năm đi học bớt cô đơn và đáng nhớ hơn nhiều.
\textit{(It makes school years less lonely and much more memorable.)}
\end{truyen}

\begin{baihoc}
Tình bạn đẹp là một trong những điều làm tuổi học trò trở nên ấm và đáng giữ.
\textit{(Beautiful friendship is one of the things that makes student life warm and worth keeping.)}
\end{baihoc}

\begin{ghinhoanh}
\textbf{Short English to remember:}

Beautiful friendship adds warmth to growing up.

\end{ghinhoanh}

\begin{tuhocnhanh}
1. Tình bạn đẹp có những nét nào? \textit{What qualities belong to beautiful friendship?}

2. Vì sao nó quý trong tuổi học trò? \textit{Why is it precious in student life?}

3. Em muốn giữ nét đẹp nào trong tình bạn của mình? \textit{What quality do you want to keep in your friendships?}
\end{tuhocnhanh}
\ngancach